\documentclass[12pt]{article}
\usepackage[a4paper, margin=2cm]{geometry}

\usepackage{fontspec}
\usepackage{polyglossia}
\setdefaultlanguage{russian}
\setmainlanguage{russian}
\setmainfont{Times New Roman}
\begin{document}

\begin{center}
\textbf{\LARGE ТЕХНИЧЕСКИЙ ОТЧЕТ}
\end{center}

\vspace{1cm}

\section*{Объект обследования}
Насосная установка (модель: \underline{\hspace{2.5cm}}, серийный номер: \underline{\hspace{2.5cm}}, тип насоса: \underline{\hspace{2.5cm}}). Расположена в помещении \underline{\hspace{3cm}}. Цель обследования — проверить работоспособность установки, зафиксировать давление в рабочем диапазоне, осмотреть узлы и оценить текущее состояние.

\section*{Цели и объём работ}
Проведение обследования в части соблюдения требований безопасности и эксплуатационной документации, выполнение измерений давления, визуальный осмотр узлов и сборки, оценка общего состояния оборудования.

\section*{Методы обследования}
\begin{itemize}
\item Визуальный осмотр кожухов, уплотнений, креплений, кабелей и состояния заземления.
\item Измерение давления в рабочей линии: входное \underline{\hspace{2.5cm}} кПа, выходное \underline{\hspace{2.5cm}} кПа.
\item Функциональная проверка: запуск/остановка, контроль шумов и вибраций.
\end{itemize}

\section*{Ход работ}
Дата обследования: \underline{\hspace{3cm}}. Применяемые методики и проведенные мероприятия отражены в приложении к данному документу.

\section*{Результаты обследования}
\begin{itemize}
\item Давление: входное = \underline{\hspace{2.5cm}} кПа; выходное = \underline{\hspace{2.5cm}} кПа; допустимый диапазон = \underline{\hspace{2.5cm}} кПа.
\item Визуальный осмотр узлов: признаков утечек, износа или повреждений не выявлено.
\item Электрическая часть: кабели без повреждений, заземление исправно, соединения плотные.
\item Механическая часть: отсутствие необычных шумов и вибраций; вращающееся оборудование функционирует нормально.
\end{itemize}

\section*{Выводы}
Оборудование работает нормально, неисправностей не обнаружено. Работы выполнены в соответствии с требованиями безопасности и нормативными документами (ГОСТы и ФНП). Рекомендации: поддерживать график планово-профилактического обслуживания и при необходимости повторить обследование по регламенту.

\section*{Приложения}
Протокол измерений и копии регламентирующих документов могут быть приложены по требованию.

\vfill

\noindent
Дата: \rule{5cm}{0.4pt} \hfill Подпись: \rule{5cm}{0.4pt}

\end{document}